\sectionheading{5.47 Exercise — Algorithm, Flowchart ও Python Program}

এই অংশে প্রতিটি Problem-এর
\textbf{Algorithm, Flowchart, Python Code এবং Output/ব্যাখ্যা}
একই Table-এর মধ্যে দেওয়া হয়েছে। ফলে একটি Problem-এর সম্পূর্ণ
Solution একসঙ্গে দেখা যাবে।

\begin{rememberbox}

\begin{center}

\textbf{Problem}
\(\longrightarrow\)
\textbf{Algorithm}
\(\longrightarrow\)
\textbf{Flowchart}
\(\longrightarrow\)
\textbf{Python Code}
\(\longrightarrow\)
\textbf{Output}

\end{center}

\end{rememberbox}

%================================================
% PART A
%================================================

\partheading{অংশ–A : Basic Input, Output ও Arithmetic}

%================================================
% 1
%================================================

\exerciseTable
{১. তিনটি সংখ্যার যোগফল ও গড় নির্ণয়}
{
১. Start।\\
২. \(a,b,c\) Input নাও।\\
৩. \(sum=a+b+c\)।\\
৪. \(average=sum/3\)।\\
৫. Sum ও Average প্রদর্শন করো।\\
৬. End।
}
{
\linearflow
{a,b,c Input}
{sum=a+b+c\\average=sum/3}
{sum, average প্রদর্শন}
}
{
a = float(input("Enter first number: "))\\
b = float(input("Enter second number: "))\\
c = float(input("Enter third number: "))\\[0.5mm]
total = a + b + c\\
average = total / 3\\
print("Sum =", total)\\
print("Average =", average)
}
{
Input:
\textenglish{\ttfamily 10, 20, 30}

\medskip

Output:

\textenglish{\ttfamily Sum = 60.0}

\textenglish{\ttfamily Average = 20.0}

\medskip

Input তিনটি Number।
Processing হলো Sum ও Average নির্ণয়।
}

%================================================
% 2
%================================================

\exerciseTable
{২. দুটি সংখ্যার যোগফল}
{
১. Start।\\
২. \(a,b\) Input নাও।\\
৩. \(sum=a+b\)।\\
৪. Sum প্রদর্শন করো।\\
৫. End।
}
{
\linearflow
{a,b Input}
{sum=a+b}
{sum প্রদর্শন}
}
{
a = int(input("Enter a: "))\\
b = int(input("Enter b: "))\\
sum = a + b\\
print("Sum =", sum)
}
{
Input:
\textenglish{\ttfamily 12, 8}

\medskip

Output:

\textenglish{\ttfamily Sum = 20}
}

%================================================
% 3
%================================================

\exerciseTable
{৩. একটি Character Input নিয়ে প্রদর্শন}
{
১. Start।\\
২. Character Input নাও।\\
৩. Character প্রদর্শন করো।\\
৪. End।
}
{
\linearflow
{Character Input}
{Character সংরক্ষণ}
{Character প্রদর্শন}
}
{
ch = input("Enter a character: ")\\
print("Character:", ch)
}
{
Input:
\textenglish{\ttfamily A}

\medskip

Output:

\textenglish{\ttfamily Character: A}
}

%================================================
% 4
%================================================

\exerciseTable
{৪. Character, String ও Number Input নিয়ে প্রদর্শন}
{
১. Start।\\
২. Character Input নাও।\\
৩. String Input নাও।\\
৪. Number Input নাও।\\
৫. সব Value প্রদর্শন করো।\\
৬. End।
}
{
\linearflow
{Character, String, Number Input}
{Valueগুলো সংরক্ষণ}
{সব Value প্রদর্শন}
}
{
ch = input("Enter a character: ")\\
text = input("Enter a string: ")\\
number = int(input("Enter a number: "))\\[0.5mm]
print("Character:", ch)\\
print("String:", text)\\
print("Number:", number)
}
{
\textenglish{\ttfamily input()}
সাধারণত String দেয়।

\medskip

Number-এর ক্ষেত্রে
\textenglish{\ttfamily int()}
ব্যবহার করা হয়েছে।
}

%================================================
% 5
%================================================

\exerciseTable
{৫. বয়সের সম্পর্ক থেকে তিনজনের মোট বয়স}
{
১. Start।\\
২. Rakib-এর বয়স \(r\) Input নাও।\\
৩. \(a=5r\)।\\
৪. \(f=a+10\)।\\
৫. \(total=r+a+f\)।\\
৬. Total প্রদর্শন করো।\\
৭. End।
}
{
\linearflow
{r Input}
{a=5r,\; f=a+10\\total=r+a+f}
{total প্রদর্শন}
}
{
r = int(input("Rakib's age: "))\\
a = 5 * r\\
f = a + 10\\
total = r + a + f\\
print("Total age =", total)
}
{
একটি Input থেকে সম্পর্ক অনুযায়ী অন্য বয়সগুলো নির্ণয় করা হয়েছে।
}

%================================================
% PART B
%================================================

\partheading{অংশ–B : Conditional Program}

%================================================
% 6
%================================================

\exerciseTable
{৬. Pass/Fail নির্ণয়}
{
১. Start।\\
২. Marks Input নাও।\\
৩. Marks \(\ge33\) কি না পরীক্ষা করো।\\
৪. সত্য হলে Pass, অন্যথায় Fail।\\
৫. End।
}
{
\binaryflow
{Marks Input}
{marks \(\ge33\)?}
{Pass}
{Fail}
}
{
marks = int(input("Enter marks: "))\\[0.5mm]
if marks >= 33:\\
\hspace*{4mm}print("Pass")\\
else:\\
\hspace*{4mm}print("Fail")
}
{
Input:
\textenglish{\ttfamily 65}

\medskip

Output:

\textenglish{\ttfamily Pass}
}

%================================================
% 7
%================================================

\exerciseTable
{৭. সংখ্যা Positive, Negative নাকি Zero}
{
১. Start।\\
২. \(n\) Input নাও।\\
৩. \(n>0\) হলে Positive।\\
৪. \(n<0\) হলে Negative।\\
৫. অন্যথায় Zero।\\
৬. End।
}
{
\linearflow
{n Input}
{n-এর Sign পরীক্ষা}
{Positive / Negative / Zero}
}
{
n = float(input("Enter a number: "))\\[0.5mm]
if n > 0:\\
\hspace*{4mm}print("Positive")\\
elif n < 0:\\
\hspace*{4mm}print("Negative")\\
else:\\
\hspace*{4mm}print("Zero")
}
{
তিনটি সম্ভাবনা থাকায়
\textenglish{\ttfamily if-elif-else}
ব্যবহার করা হয়েছে।
}

%================================================
% 8
%================================================

\exerciseTable
{৮. দুটি সংখ্যার মধ্যে বড় সংখ্যা}
{
১. Start।\\
২. \(a,b\) Input নাও।\\
৩. \(a>b\) হলে \(a\) বড়।\\
৪. \(b>a\) হলে \(b\) বড়।\\
৫. অন্যথায় দুটি সমান।\\
৬. End।
}
{
\linearflow
{a,b Input}
{দুটি সংখ্যা তুলনা}
{বড় সংখ্যা / Equal}
}
{
a = int(input("Enter a: "))\\
b = int(input("Enter b: "))\\[0.5mm]
if a > b:\\
\hspace*{4mm}print(a, "is larger")\\
elif b > a:\\
\hspace*{4mm}print(b, "is larger")\\
else:\\
\hspace*{4mm}print("Both are equal")
}
{
দুটি Number Comparison করে বড় Value নির্ণয় করা হয়েছে।
}

%================================================
% 9
%================================================

\exerciseTable
{৯. তিনটি সংখ্যার মধ্যে বৃহত্তম}
{
১. Start।\\
২. \(a,b,c\) Input নাও।\\
৩. \(a\) অন্য দুইটির চেয়ে বড়/সমান কি না দেখো।\\
৪. না হলে \(b\) পরীক্ষা করো।\\
৫. অন্যথায় \(c\) বৃহত্তম।\\
৬. Largest প্রদর্শন করো।\\
৭. End।
}
{
\linearflow
{a,b,c Input}
{তিনটি সংখ্যা তুলনা}
{largest প্রদর্শন}
}
{
a = int(input("Enter a: "))\\
b = int(input("Enter b: "))\\
c = int(input("Enter c: "))\\[0.5mm]
if a >= b and a >= c:\\
\hspace*{4mm}largest = a\\
elif b >= a and b >= c:\\
\hspace*{4mm}largest = b\\
else:\\
\hspace*{4mm}largest = c\\
print("Largest =", largest)
}
{
Input:
\textenglish{\ttfamily 10,25,15}

\medskip

Output:

\textenglish{\ttfamily Largest = 25}
}

%================================================
% 10
%================================================

\exerciseTable
{১০. তিনটি সংখ্যার মধ্যে ক্ষুদ্রতম}
{
১. Start।\\
২. \(a,b,c\) Input নাও।\\
৩. তিনটি Value তুলনা করো।\\
৪. Smallest নির্ধারণ করো।\\
৫. Result প্রদর্শন করো।\\
৬. End।
}
{
\linearflow
{a,b,c Input}
{তিনটি সংখ্যা তুলনা}
{smallest প্রদর্শন}
}
{
a = int(input("Enter a: "))\\
b = int(input("Enter b: "))\\
c = int(input("Enter c: "))\\[0.5mm]
if a <= b and a <= c:\\
\hspace*{4mm}smallest = a\\
elif b <= a and b <= c:\\
\hspace*{4mm}smallest = b\\
else:\\
\hspace*{4mm}smallest = c\\
print("Smallest =", smallest)
}
{
Largest-এর বিপরীত Comparison Logic ব্যবহার করা হয়েছে।
}

%================================================
% 11
%================================================

\exerciseTable
{১১. সংখ্যা জোড় না বিজোড়}
{
১. Start।\\
২. \(n\) Input নাও।\\
৩. \(n\bmod2=0\) কি না পরীক্ষা করো।\\
৪. সত্য হলে Even, অন্যথায় Odd।\\
৫. End।
}
{
\binaryflow
{n Input}
{n \% 2 == 0 ?}
{Even}
{Odd}
}
{
n = int(input("Enter n: "))\\[0.5mm]
if n \% 2 == 0:\\
\hspace*{4mm}print("Even")\\
else:\\
\hspace*{4mm}print("Odd")
}
{
\textenglish{\ttfamily \%}
দিয়ে ভাগশেষ পাওয়া যায়।

\medskip

ভাগশেষ 0 হলে সংখ্যাটি Even।
}

%================================================
% 12
%================================================

\exerciseTable
{১২. Capital, Small নাকি অন্য Character}
{
১. Start।\\
২. Character Input নাও।\\
৩. `A-Z' Range পরীক্ষা করো।\\
৪. না হলে `a-z' Range পরীক্ষা করো।\\
৫. Type প্রদর্শন করো।\\
৬. End।
}
{
\linearflow
{Character Input}
{Capital / Small Range পরীক্ষা}
{Character Type প্রদর্শন}
}
{
ch = input("Enter character: ")\\[0.5mm]
if 'A' <= ch <= 'Z':\\
\hspace*{4mm}print("Capital letter")\\
elif 'a' <= ch <= 'z':\\
\hspace*{4mm}print("Small letter")\\
else:\\
\hspace*{4mm}print("Other character")
}
{
Python-এ Character Range সরাসরি Compare করা যায়।
}

%================================================
% 13
%================================================

\exerciseTable
{১৩. Vowel নাকি Consonant}
{
১. Start।\\
২. Letter Input নাও।\\
৩. Letterটি a,e,i,o,u-এর কোনোটি কি না পরীক্ষা করো।\\
৪. Vowel বা Consonant প্রদর্শন করো।\\
৫. End।
}
{
\binaryflow
{Letter Input}
{Vowel list-এ আছে?}
{Vowel}
{Consonant}
}
{
ch = input("Enter a letter: ")\\[0.5mm]
if ch.lower() in "aeiou":\\
\hspace*{4mm}print("Vowel")\\
else:\\
\hspace*{4mm}print("Consonant")
}
{
\textenglish{\ttfamily lower()}
ব্যবহার করায় Capital ও Small উভয় Letter কাজ করবে।
}

%================================================
% 14
%================================================

\exerciseTable
{১৪. বয়স অনুযায়ী শ্রেণি নির্ধারণ}
{
১. Start।\\
২. Age Input নাও।\\
৩. Age Range পরীক্ষা করো।\\
৪. Category নির্ধারণ করো।\\
৫. Result প্রদর্শন করো।\\
৬. End।
}
{
\linearflow
{Age Input}
{Age Range পরীক্ষা}
{Category প্রদর্শন}
}
{
age = int(input("Enter age: "))\\[0.5mm]
if age <= 12:\\
\hspace*{4mm}print("Child")\\
elif age <= 19:\\
\hspace*{4mm}print("Teenager")\\
elif age <= 40:\\
\hspace*{4mm}print("Young")\\
else:\\
\hspace*{4mm}print("Other")
}
{
Rangeগুলো ক্রমানুসারে পরীক্ষা করা হয়েছে।
}

%================================================
% 15
%================================================

\exerciseTable
{১৫. ৫ ও ৭ দ্বারা বিভাজ্যতা পরীক্ষা}
{
১. Start।\\
২. \(n\) Input নাও।\\
৩. ৫ ও ৭ উভয় দ্বারা বিভাজ্য কি না দেখো।\\
৪. শুধু ৫ বা শুধু ৭ দ্বারা বিভাজ্য কি না দেখো।\\
৫. নির্ধারিত Message প্রদর্শন করো।\\
৬. End।
}
{
\linearflow
{n Input}
{5 ও 7 দ্বারা divisibility পরীক্ষা}
{Message প্রদর্শন}
}
{
n = int(input("Enter n: "))\\[0.5mm]
if n \% 5 == 0 and n \% 7 == 0:\\
\hspace*{4mm}print("Good")\\
elif n \% 5 == 0:\\
\hspace*{4mm}print("Flower")\\
elif n \% 7 == 0:\\
\hspace*{4mm}print("River")\\
else:\\
\hspace*{4mm}print("Try again")
}
{
উভয় দ্বারা বিভাজ্যতার Condition প্রথমে পরীক্ষা করতে হবে।
}

%================================================
% 16
%================================================

\exerciseTable
{১৬. Roll অনুযায়ী Group নির্ধারণ}
{
১. Start।\\
২. Roll Input নাও।\\
৩. 1--30 হলে A।\\
৪. 31--60 হলে B।\\
৫. 61--100 হলে C।\\
৬. অন্যথায় Invalid।\\
৭. End।
}
{
\linearflow
{Roll Input}
{Roll Range পরীক্ষা}
{Group প্রদর্শন}
}
{
roll = int(input("Enter roll: "))\\[0.5mm]
if 1 <= roll <= 30:\\
\hspace*{4mm}print("Group A")\\
elif roll <= 60:\\
\hspace*{4mm}print("Group B")\\
elif roll <= 100:\\
\hspace*{4mm}print("Group C")\\
else:\\
\hspace*{4mm}print("Invalid")
}
{
এটি Range-based Conditional-এর উদাহরণ।
}

%================================================
% 17
%================================================

\exerciseTable
{১৭. Marks থেকে Grade নির্ণয়}
{
১. Start।\\
২. Marks Input নাও।\\
৩. সর্বোচ্চ Range থেকে Condition পরীক্ষা করো।\\
৪. Grade নির্ধারণ করো।\\
৫. Grade প্রদর্শন করো।\\
৬. End।
}
{
\linearflow
{Marks Input}
{Marks Range পরীক্ষা}
{Grade প্রদর্শন}
}
{
marks = int(input("Enter marks: "))\\[0.5mm]
if marks >= 80:\\
\hspace*{4mm}grade = "A+"\\
elif marks >= 70:\\
\hspace*{4mm}grade = "A"\\
elif marks >= 60:\\
\hspace*{4mm}grade = "A-"\\
elif marks >= 50:\\
\hspace*{4mm}grade = "B"\\
elif marks >= 40:\\
\hspace*{4mm}grade = "C"\\
elif marks >= 33:\\
\hspace*{4mm}grade = "D"\\
else:\\
\hspace*{4mm}grade = "F"\\
print("Grade =", grade)
}
{
Condition সর্বোচ্চ Marks থেকে নিচে নামানো হয়েছে।
}

%================================================
% PART C
%================================================

\partheading{অংশ–C : Loop ও Series ভিত্তিক Program}

%================================================
% 18
%================================================

\exerciseTable
{১৮. কোনো Text নির্দিষ্ট সংখ্যকবার Print}
{
১. Start।\\
২. Text ও \(n\) Input নাও।\\
৩. Loop \(n\) বার চালাও।\\
৪. প্রতিবার Text Print করো।\\
৫. End।
}
{
\linearflow
{Text, n Input}
{Loop n বার}
{Text Print}
}
{
text = input("Enter text: ")\\
n = int(input("How many times: "))\\
for i in range(n):\\
\hspace*{4mm}print(text)
}
{
\textenglish{\ttfamily range(n)}
Loop-কে \(n\) বার চালায়।
}

%================================================
% 19
%================================================

\exerciseTable
{১৯. ১ থেকে ১০ পর্যন্ত সংখ্যা প্রদর্শন}
{
১. Start।\\
২. 1 থেকে 10 পর্যন্ত Loop চালাও।\\
৩. প্রতিটি Value Print করো।\\
৪. End।
}
{
\linearflow
{i=1}
{i=1 থেকে 10 পর্যন্ত Loop}
{i প্রদর্শন}
}
{
for i in range(1, 11):\\
\hspace*{4mm}print(i, end=" ")
}
{
Output:

\textenglish{\ttfamily 1 2 3 4 5 6 7 8 9 10}
}

%================================================
% 20
%================================================

\exerciseTable
{২০. ১০ থেকে ১ উল্টোভাবে প্রদর্শন}
{
১. Start।\\
২. \(i=10\) নাও।\\
৩. Value Print করো।\\
৪. \(i=i-1\)।\\
৫. 1 পর্যন্ত Repeat করো।\\
৬. End।
}
{
\linearflow
{i=10}
{প্রতি ধাপে i এক কমাও}
{10 থেকে 1 Print}
}
{
for i in range(10, 0, -1):\\
\hspace*{4mm}print(i, end=" ")
}
{
Negative Step ব্যবহার করা হয়েছে।
}

%================================================
% 21
%================================================

\exerciseTable
{২১. ১ থেকে n পর্যন্ত Odd Number}
{
১. Start।\\
২. \(n\) Input নাও।\\
৩. 1 থেকে শুরু করো।\\
৪. Step 2 করে \(n\) পর্যন্ত Print করো।\\
৫. End।
}
{
\linearflow
{n Input}
{1 থেকে n, Step=2}
{Odd Number Print}
}
{
n = int(input("Enter n: "))\\
for i in range(1, n + 1, 2):\\
\hspace*{4mm}print(i, end=" ")
}
{
Step 2 হওয়ায় শুধু Odd Number পাওয়া যায়।
}

%================================================
% 22
%================================================

\exerciseTable
{২২. ১ থেকে n পর্যন্ত Even Number}
{
১. Start।\\
২. \(n\) Input নাও।\\
৩. 2 থেকে শুরু করো।\\
৪. Step 2 করে \(n\) পর্যন্ত Print করো।\\
৫. End।
}
{
\linearflow
{n Input}
{2 থেকে n, Step=2}
{Even Number Print}
}
{
n = int(input("Enter n: "))\\
for i in range(2, n + 1, 2):\\
\hspace*{4mm}print(i, end=" ")
}
{
Loop সরাসরি Even Value তৈরি করছে।
}

%================================================
% 23
%================================================

\exerciseTable
{২৩. ১০ থেকে ১০০ পর্যন্ত ১০-এর ব্যবধানে সংখ্যা}
{
১. Start।\\
২. 10 থেকে শুরু করো।\\
৩. Step 10 করে 100 পর্যন্ত Value নাও।\\
৪. প্রতিটি Value Print করো।\\
৫. End।
}
{
\linearflow
{i=10}
{i=10,20,\dots,100}
{i প্রদর্শন}
}
{
for i in range(10, 101, 10):\\
\hspace*{4mm}print(i, end=" ")
}
{
\textenglish{\ttfamily range(start, stop, step)}
ব্যবহার করা হয়েছে।
}

%================================================
% 24
%================================================

\exerciseTable
{২৪. ছোট হাতের a--z প্রদর্শন}
{
১. Start।\\
২. a-এর Character Code থেকে শুরু করো।\\
৩. z পর্যন্ত প্রতিটি Code Traverse করো।\\
৪. Character Print করো।\\
৫. End।
}
{
\linearflow
{a থেকে z}
{Character Code Traverse}
{a--z প্রদর্শন}
}
{
for code in range(ord('a'), ord('z') + 1):\\
\hspace*{4mm}print(chr(code), end=" ")
}
{
\textenglish{\ttfamily ord()}
Character-এর Code দেয়।

\medskip

\textenglish{\ttfamily chr()}
Code থেকে Character দেয়।
}

%================================================
% 25
%================================================

\exerciseTable
{২৫. বড় হাতের A--Z প্রদর্শন}
{
১. Start।\\
২. A-এর Character Code থেকে শুরু করো।\\
৩. Z পর্যন্ত Traverse করো।\\
৪. Character Print করো।\\
৫. End।
}
{
\linearflow
{A থেকে Z}
{Character Code Traverse}
{A--Z প্রদর্শন}
}
{
for code in range(ord('A'), ord('Z') + 1):\\
\hspace*{4mm}print(chr(code), end=" ")
}
{
Capital Alphabet ধারাবাহিকভাবে Print হবে।
}

%================================================
% 26
%================================================

\exerciseTable
{২৬. ১ থেকে ১০০-এর মধ্যে ৭ দ্বারা বিভাজ্য সংখ্যা}
{
১. Start।\\
২. 1 থেকে 100 পর্যন্ত Loop চালাও।\\
৩. \(i\bmod7=0\) পরীক্ষা করো।\\
৪. সত্য হলে \(i\) Print করো।\\
৫. End।
}
{
\linearflow
{1 থেকে 100}
{i \% 7 == 0 পরীক্ষা}
{সত্য হলে Print}
}
{
for i in range(1, 101):\\
\hspace*{4mm}if i \% 7 == 0:\\
\hspace*{8mm}print(i, end=" ")
}
{
Condition True হওয়া Value-গুলোই Print হয়।
}

%================================================
% 27
%================================================

\exerciseTable
{২৭. একটি সংখ্যার সব Divisor}
{
১. Start।\\
২. \(n\) Input নাও।\\
৩. 1 থেকে \(n\) পর্যন্ত Loop চালাও।\\
৪. \(n\bmod i=0\) হলে \(i\) Print করো।\\
৫. End।
}
{
\linearflow
{n Input}
{1..n দিয়ে divisibility পরীক্ষা}
{Divisor Print}
}
{
n = int(input("Enter n: "))\\
for i in range(1, n + 1):\\
\hspace*{4mm}if n \% i == 0:\\
\hspace*{8mm}print(i, end=" ")
}
{
যে Number দিয়ে নিঃশেষে ভাগ যায় সেটিই Divisor।
}

%================================================
% 28
%================================================

\exerciseTable
{২৮. ১ থেকে ৫ পর্যন্ত সংখ্যার Square}
{
১. Start।\\
২. 1 থেকে 5 পর্যন্ত Loop চালাও।\\
৩. প্রতিটি Value-এর Square নির্ণয় করো।\\
৪. Result Print করো।\\
৫. End।
}
{
\linearflow
{i=1..5}
{square=i*i}
{i ও square Print}
}
{
for i in range(1, 6):\\
\hspace*{4mm}print(i, "square =", i * i)
}
{
প্রতিটি Number-এর Square প্রদর্শিত হবে।
}

%================================================
% 29
%================================================

\exerciseTable
{২৯. ১ থেকে ১০ পর্যন্ত সংখ্যার Cube-এর যোগফল}
{
১. Start।\\
২. \(sum=0\)।\\
৩. 1 থেকে 10 পর্যন্ত Loop চালাও।\\
৪. \(sum=sum+i^3\)।\\
৫. Sum প্রদর্শন করো।\\
৬. End।
}
{
\linearflow
{i=1..10}
{sum=sum+i\(^3\)}
{sum প্রদর্শন}
}
{
total = 0\\
for i in range(1, 11):\\
\hspace*{4mm}total += i ** 3\\
print("Sum =", total)
}
{
\textenglish{\ttfamily **}
হলো Python-এর Power Operator।
}

%================================================
% 30
%================================================

\exerciseTable
{৩০. ১ থেকে n পর্যন্ত সংখ্যার যোগফল}
{
১. Start।\\
২. \(n\) Input নাও।\\
৩. \(sum=0\)।\\
৪. 1 থেকে \(n\) পর্যন্ত Value যোগ করো।\\
৫. Sum প্রদর্শন করো।\\
৬. End।
}
{
\linearflow
{n Input}
{1..n Accumulate}
{sum প্রদর্শন}
}
{
n = int(input("Enter n: "))\\
total = 0\\
for i in range(1, n + 1):\\
\hspace*{4mm}total += i\\
print("Sum =", total)
}
{
\textenglish{\ttfamily total}
Accumulator হিসেবে কাজ করছে।
}

%================================================
% 31
%================================================

\exerciseTable
{৩১. ১ থেকে n পর্যন্ত Even সংখ্যার যোগফল}
{
১. Start।\\
২. \(n\) Input নাও।\\
৩. Even Value তৈরি করো।\\
৪. প্রতিটি Even Value Sum-এ যোগ করো।\\
৫. Result Print করো।\\
৬. End।
}
{
\linearflow
{n Input}
{Even Value Accumulate}
{Even Sum Print}
}
{
n = int(input("Enter n: "))\\
total = 0\\
for i in range(2, n + 1, 2):\\
\hspace*{4mm}total += i\\
print("Even sum =", total)
}
{
Step 2 ব্যবহার করে সরাসরি Even Value নেওয়া হয়েছে।
}

%================================================
% 32
%================================================

\exerciseTable
{৩২. ১ থেকে n পর্যন্ত Odd সংখ্যার যোগফল}
{
১. Start।\\
২. \(n\) Input নাও।\\
৩. Odd Value তৈরি করো।\\
৪. প্রতিটি Odd Value Sum-এ যোগ করো।\\
৫. Result Print করো।\\
৬. End।
}
{
\linearflow
{n Input}
{Odd Value Accumulate}
{Odd Sum Print}
}
{
n = int(input("Enter n: "))\\
total = 0\\
for i in range(1, n + 1, 2):\\
\hspace*{4mm}total += i\\
print("Odd sum =", total)
}
{
Step 2 দিয়ে Odd Series পাওয়া যায়।
}

%================================================
% 33
%================================================

\exerciseTable
{৩৩. \(10+9+8+\cdots+Y\)}
{
১. Start।\\
২. \(Y\) Input নাও।\\
৩. \(i=10\), \(sum=0\) নাও।\\
৪. \(Y\) পর্যন্ত Descending Value যোগ করো।\\
৫. Sum প্রদর্শন করো।\\
৬. End।
}
{
\linearflow
{Y Input}
{10 থেকে Y পর্যন্ত Descending Sum}
{sum প্রদর্শন}
}
{
y = int(input("Enter Y: "))\\
total = 0\\
for i in range(10, y - 1, -1):\\
\hspace*{4mm}total += i\\
print("Sum =", total)
}
{
Negative Step ব্যবহার করা হয়েছে।
}

%================================================
% 34
%================================================

\exerciseTable
{৩৪. \(5+10+15+\cdots+n\)}
{
১. Start।\\
২. \(n\) Input নাও।\\
৩. 5 থেকে Step 5 করে Value তৈরি করো।\\
৪. প্রতিটি Value Sum-এ যোগ করো।\\
৫. Result Print করো।\\
৬. End।
}
{
\linearflow
{n Input}
{5,10,15,\dots,n যোগ}
{sum প্রদর্শন}
}
{
n = int(input("Enter n: "))\\
total = 0\\
for i in range(5, n + 1, 5):\\
\hspace*{4mm}total += i\\
print(total)
}
{
Step 5 প্রতিটি পরবর্তী Term তৈরি করছে।
}

%================================================
% 35
%================================================

\exerciseTable
{৩৫. \(1^2+2^2+3^2+\cdots+n^2\)}
{
১. Start।\\
২. \(n\) Input নাও।\\
৩. 1 থেকে \(n\) পর্যন্ত Loop চালাও।\\
৪. \(i^2\) Sum-এ যোগ করো।\\
৫. Result Print করো।\\
৬. End।
}
{
\linearflow
{n Input}
{sum=sum+i\(^2\)}
{sum প্রদর্শন}
}
{
n = int(input("Enter n: "))\\
total = 0\\
for i in range(1, n + 1):\\
\hspace*{4mm}total += i ** 2\\
print(total)
}
{
প্রতিটি Term হলো \(i^2\)।
}

%================================================
% 36
%================================================

\exerciseTable
{৩৬. \(1^2+3^2+5^2+\cdots+17^2\)}
{
১. Start।\\
২. 1 থেকে 17 পর্যন্ত Odd Value নাও।\\
৩. প্রতিটি Value-এর Square নির্ণয় করো।\\
৪. Squareগুলো যোগ করো।\\
৫. Result Print করো।\\
৬. End।
}
{
\linearflow
{Odd 1..17}
{sum=sum+i\(^2\)}
{sum প্রদর্শন}
}
{
total = 0\\
for i in range(1, 18, 2):\\
\hspace*{4mm}total += i ** 2\\
print(total)
}
{
শুধু Odd Number-এর Square যোগ করা হয়েছে।
}

%================================================
% 37
%================================================

\exerciseTable
{৩৭. \(2^2+4^2+6^2+\cdots+n^2\)}
{
১. Start।\\
২. \(n\) Input নাও।\\
৩. Even Value তৈরি করো।\\
৪. প্রতিটি Value-এর Square যোগ করো।\\
৫. Result Print করো।\\
৬. End।
}
{
\linearflow
{n Input}
{Even i-এর Square যোগ}
{sum প্রদর্শন}
}
{
n = int(input("Enter n: "))\\
total = 0\\
for i in range(2, n + 1, 2):\\
\hspace*{4mm}total += i ** 2\\
print(total)
}
{
শুধু Even base ব্যবহার করা হয়েছে।
}

%================================================
% 38
%================================================

\exerciseTable
{৩৮. \(1^5+2^5+\cdots+n^5\)}
{
১. Start।\\
২. \(n\) Input নাও।\\
৩. 1 থেকে \(n\) পর্যন্ত Loop চালাও।\\
৪. \(i^5\) যোগ করো।\\
৫. Result Print করো।\\
৬. End।
}
{
\linearflow
{n Input}
{sum=sum+i\(^5\)}
{sum প্রদর্শন}
}
{
n = int(input("Enter n: "))\\
total = 0\\
for i in range(1, n + 1):\\
\hspace*{4mm}total += i ** 5\\
print(total)
}
{
\textenglish{\ttfamily i ** 5}
মানে \(i^5\)।
}

%================================================
% 39
%================================================

\exerciseTable
{৩৯. \(3^6+5^6+7^6+\cdots+n^6\)}
{
১. Start।\\
২. \(n\) Input নাও।\\
৩. 3 থেকে Odd Value তৈরি করো।\\
৪. প্রতিটি Value-এর 6th Power যোগ করো।\\
৫. Result Print করো।\\
৬. End।
}
{
\linearflow
{n Input}
{Odd Term-এর 6th Power যোগ}
{sum প্রদর্শন}
}
{
n = int(input("Enter n: "))\\
total = 0\\
for i in range(3, n + 1, 2):\\
\hspace*{4mm}total += i ** 6\\
print(total)
}
{
3, 5, 7, \dots Series ব্যবহৃত হয়েছে।
}

%================================================
% 40
%================================================

\exerciseTable
{৪০. \(2^4+4^4+6^4+\cdots+n^4\)}
{
১. Start।\\
২. \(n\) Input নাও।\\
৩. Even Value তৈরি করো।\\
৪. প্রতিটি Value-এর 4th Power যোগ করো।\\
৫. Result Print করো।\\
৬. End।
}
{
\linearflow
{n Input}
{Even Term-এর 4th Power যোগ}
{sum প্রদর্শন}
}
{
n = int(input("Enter n: "))\\
total = 0\\
for i in range(2, n + 1, 2):\\
\hspace*{4mm}total += i ** 4\\
print(total)
}
{
Even Terms-এর Fourth Power যোগ করা হয়েছে।
}

%================================================
% 41
%================================================

\exerciseTable
{৪১. \(1+\frac{1}{2^2}+\frac{1}{3^2}+\cdots+\frac{1}{n^2}\)}
{
১. Start।\\
২. \(n\) Input নাও।\\
৩. 1 থেকে \(n\) পর্যন্ত Loop চালাও।\\
৪. \(1/i^2\) যোগ করো।\\
৫. Result Print করো।\\
৬. End।
}
{
\linearflow
{n Input}
{sum=sum+1/i\(^2\)}
{sum প্রদর্শন}
}
{
n = int(input("Enter n: "))\\
total = 0.0\\
for i in range(1, n + 1):\\
\hspace*{4mm}total += 1 / (i ** 2)\\
print(total)
}
{
Fraction থাকায় Result সাধারণত Float হবে।
}

%================================================
% 42
%================================================

\exerciseTable
{৪২. \(1\times2+2\times3+3\times4+\cdots+n(n+1)\)}
{
১. Start।\\
২. \(n\) Input নাও।\\
৩. 1 থেকে \(n\) পর্যন্ত Loop চালাও।\\
৪. \(i(i+1)\) যোগ করো।\\
৫. Result Print করো।\\
৬. End।
}
{
\linearflow
{n Input}
{sum=sum+i(i+1)}
{sum প্রদর্শন}
}
{
n = int(input("Enter n: "))\\
total = 0\\
for i in range(1, n + 1):\\
\hspace*{4mm}total += i * (i + 1)\\
print(total)
}
{
প্রতিটি Term-এ দুটি ধারাবাহিক Integer গুণ করা হয়েছে।
}

%================================================
% 43
%================================================

\exerciseTable
{৪৩. Factorial নির্ণয়}
{
১. Start।\\
২. \(n\) Input নাও।\\
৩. \(fact=1\)।\\
৪. 1 থেকে \(n\) পর্যন্ত সব Value দ্বারা fact গুণ করো।\\
৫. Result প্রদর্শন করো।\\
৬. End।
}
{
\linearflow
{n Input}
{fact=1; 1..n Multiply}
{fact প্রদর্শন}
}
{
n = int(input("Enter n: "))\\
fact = 1\\
for i in range(1, n + 1):\\
\hspace*{4mm}fact *= i\\
print("Factorial =", fact)
}
{
Input 5 হলে—

\textenglish{\ttfamily Factorial = 120}

\medskip

Multiplication-এর Initial Value 1।
}

%================================================
% 44
%================================================

\exerciseTable
{৪৪. Geometric Series: \(1+3+9+27+\cdots+N\)}
{
১. Start।\\
২. \(N\) Input নাও।\\
৩. term=1, sum=0।\\
৪. term \(\le N\) পর্যন্ত যোগ করো।\\
৫. term-কে 3 দিয়ে গুণ করো।\\
৬. Sum প্রদর্শন করো।\\
৭. End।
}
{
\linearflow
{N Input}
{sum+=term; term*=3}
{sum প্রদর্শন}
}
{
n = int(input("Enter last term: "))\\
term = 1\\
total = 0\\
while term <= n:\\
\hspace*{4mm}total += term\\
\hspace*{4mm}term *= 3\\
print("Sum =", total)
}
{
প্রতিটি নতুন Term আগের Term-এর 3 গুণ।
}

%================================================
% 45
%================================================

\exerciseTable
{৪৫. Increasing Number Pattern}
{
১. Start।\\
২. Outer Loop Row নিয়ন্ত্রণ করবে।\\
৩. Inner Loop 1 থেকে Row Number পর্যন্ত চলবে।\\
৪. প্রতিটি Row শেষে New Line।\\
৫. End।
}
{
\linearflow
{Row 1..5}
{প্রতি Row-এ 1..i Print}
{Pattern প্রদর্শন}
}
{
for i in range(1, 6):\\
\hspace*{4mm}for j in range(1, i + 1):\\
\hspace*{8mm}print(j, end=" ")\\
\hspace*{4mm}print()
}
{
\begin{english}
\ttfamily
1\\
1 2\\
1 2 3\\
1 2 3 4\\
1 2 3 4 5
\end{english}
}

%================================================
% 46
%================================================

\exerciseTable
{৪৬. ২ থেকে n পর্যন্ত Prime Number}
{
১. Start।\\
২. \(n\) Input নাও।\\
৩. 2 থেকে \(n\) পর্যন্ত প্রতিটি Number পরীক্ষা করো।\\
৪. Divisor পাওয়া গেলে Prime নয়।\\
৫. Divisor না পেলে Prime Print করো।\\
৬. End।
}
{
\linearflow
{n Input}
{প্রতিটি Number-এর Divisor পরীক্ষা}
{Prime হলে Print}
}
{
n = int(input("Enter highest range: "))\\
for num in range(2, n + 1):\\
\hspace*{4mm}is\_prime = True\\
\hspace*{4mm}for d in range(2, num):\\
\hspace*{8mm}if num \% d == 0:\\
\hspace*{12mm}is\_prime = False\\
\hspace*{12mm}break\\
\hspace*{4mm}if is\_prime:\\
\hspace*{8mm}print(num, end=" ")
}
{
Divisor পাওয়া গেলে
\textenglish{\ttfamily break}
দিয়ে Inner Loop বন্ধ করা হয়েছে।
}

%================================================
% 47
%================================================

\exerciseTable
{৪৭. প্রতি ৭ দিন পরপর জানুয়ারি মাসের Date}
{
১. Start।\\
২. Day = 1 নাও।\\
৩. Step 7 করে 31 পর্যন্ত Date তৈরি করো।\\
৪. প্রতিটি Date Print করো।\\
৫. End।
}
{
\linearflow
{day=1}
{Step=7,\; day \(\le31\)}
{Date প্রদর্শন}
}
{
for day in range(1, 32, 7):\\
\hspace*{4mm}print(day, "January")
}
{
Output:

\textenglish{\ttfamily 1, 8, 15, 22, 29}
}

%================================================
% 48
%================================================

\exerciseTable
{৪৮. ১০–৫০-এর মধ্যে ১৫ ও ২৫ বাদ দিয়ে Print}
{
১. Start।\\
২. 10 থেকে 50 পর্যন্ত Loop চালাও।\\
৩. 15 বা 25 হলে Skip করো।\\
৪. অন্য Value Print করো।\\
৫. End।
}
{
\linearflow
{i=10..50}
{i=15 বা 25 হলে continue}
{অন্য Value Print}
}
{
for i in range(10, 51):\\
\hspace*{4mm}if i == 15 or i == 25:\\
\hspace*{8mm}continue\\
\hspace*{4mm}print(i, end=" ")
}
{
15 এবং 25 Output-এ থাকবে না।
}

%================================================
% 49
%================================================

\exerciseTable
{৪৯. Positive Number না পাওয়া পর্যন্ত Input নেওয়া}
{
১. Start।\\
২. Number Input নাও।\\
৩. \(n\le0\) হলে আবার Input নাও।\\
৪. Positive হলে Loop বন্ধ করো।\\
৫. Number প্রদর্শন করো।\\
৬. End।
}
{
\linearflow
{Number Input}
{Positive না হওয়া পর্যন্ত Repeat}
{Positive Number প্রদর্শন}
}
{
while True:\\
\hspace*{4mm}n = int(input("Enter positive: "))\\
\hspace*{4mm}if n <= 0:\\
\hspace*{8mm}continue\\
\hspace*{4mm}break\\
print(n)
}
{
Positive Input পাওয়া গেলে
\textenglish{\ttfamily break}
Loop বন্ধ করে।
}

%================================================
% PART D
%================================================

\partheading{অংশ–D : List, String ও Search}

%================================================
% 50
%================================================

\exerciseTable
{৫০. List-এর সব Element প্রদর্শন}
{
১. Start।\\
২. List তৈরি করো।\\
৩. প্রতিটি Element Traverse করো।\\
৪. Element Print করো।\\
৫. End।
}
{
\linearflow
{List}
{প্রতিটি Element Traverse}
{Element Print}
}
{
numbers = [65, 78, 84]\\
for x in numbers:\\
\hspace*{4mm}print(x)
}
{
List-এর প্রতিটি Item একবার করে পাওয়া যায়।
}

%================================================
% 51
%================================================

\exerciseTable
{৫১. List-এর সব সংখ্যার যোগফল}
{
১. Start।\\
২. List তৈরি করো।\\
৩. \(total=0\) নাও।\\
৪. প্রতিটি Item total-এ যোগ করো।\\
৫. Result Print করো।\\
৬. End।
}
{
\linearflow
{List}
{সব Item Accumulate}
{total প্রদর্শন}
}
{
numbers = [12, 32, 34, 54, 65, 16, 21]\\
total = 0\\
for x in numbers:\\
\hspace*{4mm}total += x\\
print("Sum =", total)
}
{
\textenglish{\ttfamily total}
Accumulator।
}

%================================================
% 52
%================================================

\exerciseTable
{৫২. User Input থেকে List তৈরি করে Average}
{
১. Start।\\
২. \(n\) Input নাও।\\
৩. \(n\)টি Number List-এ রাখো।\\
৪. Sum-কে \(n\) দিয়ে ভাগ করো।\\
৫. Average Print করো।\\
৬. End।
}
{
\linearflow
{n ও nটি Number Input}
{average=sum(numbers)/n}
{average প্রদর্শন}
}
{
n = int(input("How many numbers: "))\\
numbers = []\\
for i in range(n):\\
\hspace*{4mm}x = float(input("Enter number: "))\\
\hspace*{4mm}numbers.append(x)\\
average = sum(numbers) / n\\
print("Average =", average)
}
{
প্রতিটি Input
\textenglish{\ttfamily append()}
দিয়ে List-এ রাখা হয়েছে।
}

%================================================
% 53
%================================================

\exerciseTable
{৫৩. String-এর Length Loop দিয়ে নির্ণয়}
{
১. Start।\\
২. String Input নাও।\\
৩. count = 0 নাও।\\
৪. প্রতিটি Character-এর জন্য count এক বাড়াও।\\
৫. Length Print করো।\\
৬. End।
}
{
\linearflow
{String Input}
{প্রতিটি Character Count}
{Length প্রদর্শন}
}
{
text = input("Enter a string: ")\\
count = 0\\
for ch in text:\\
\hspace*{4mm}count += 1\\
print("Length =", count)
}
{
Built-in
\textenglish{\ttfamily len()}
ব্যবহার না করে Counting Logic দেখানো হয়েছে।
}

%================================================
% 54
%================================================

\exerciseTable
{৫৪. List-এ নির্দিষ্ট সংখ্যা Search}
{
১. Start।\\
২. List ও Key নাও।\\
৩. প্রতিটি Item-এর সঙ্গে Key তুলনা করো।\\
৪. Match হলে Found এবং Loop বন্ধ।\\
৫. না হলে Not Found।\\
৬. End।
}
{
\linearflow
{List ও Key}
{প্রতিটি Item-এর সঙ্গে Key তুলনা}
{Found / Not Found}
}
{
numbers = [3, 5, 6, 9, 12, 21, 27, 31]\\
key = int(input("Search: "))\\
found = False\\
for x in numbers:\\
\hspace*{4mm}if x == key:\\
\hspace*{8mm}found = True\\
\hspace*{8mm}break\\
if found:\\
\hspace*{4mm}print("Found")\\
else:\\
\hspace*{4mm}print("Not found")
}
{
প্রতিটি Item একে একে পরীক্ষা করা হয়েছে।

\medskip

এটি Linear Search।
}

%================================================
% 55
%================================================

\exerciseTable
{৫৫. List-এর বৃহত্তম সংখ্যা}
{
১. Start।\\
২. প্রথম Item-কে largest ধরো।\\
৩. বাকি Item Traverse করো।\\
৪. বড় Value পেলে largest Update করো।\\
৫. Result Print করো।\\
৬. End।
}
{
\linearflow
{List}
{Compare ও largest Update}
{largest প্রদর্শন}
}
{
numbers = [3, 5, 9, 4, 13]\\
largest = numbers[0]\\
for x in numbers:\\
\hspace*{4mm}if x > largest:\\
\hspace*{8mm}largest = x\\
print("Largest =", largest)
}
{
Output:

\textenglish{\ttfamily Largest = 13}
}

%================================================
% 56
%================================================

\exerciseTable
{৫৬. দুটি String তুলনা}
{
১. Start।\\
২. দুটি String Input নাও।\\
৩. Equality পরীক্ষা করো।\\
৪. Equal না হলে Lexicographic Comparison করো।\\
৫. Result Print করো।\\
৬. End।
}
{
\linearflow
{s1,s2 Input}
{String Comparison}
{Result প্রদর্শন}
}
{
s1 = input("First string: ")\\
s2 = input("Second string: ")\\
if s1 == s2:\\
\hspace*{4mm}print("Equal")\\
elif s1 > s2:\\
\hspace*{4mm}print(s1, "comes after", s2)\\
else:\\
\hspace*{4mm}print(s1, "comes before", s2)
}
{
Python-এ String সরাসরি Compare করা যায়।
}

%================================================
% PART E
%================================================

\partheading{অংশ–E : Function ভিত্তিক Program}

%================================================
% 57
%================================================

\exerciseTable
{৫৭. Function দিয়ে তিনটি সংখ্যার যোগফল}
{
১. Start।\\
২. তিনটি Number Input নাও।\\
৩. \textenglish{\ttfamily add()} Function Call করো।\\
৪. Function Sum Return করবে।\\
৫. Result প্রদর্শন করো।\\
৬. End।
}
{
\linearflow
{x,y,z Input}
{add(x,y,z) Call}
{Returned Sum Print}
}
{
def add(a, b, c):\\
\hspace*{4mm}return a + b + c\\
x = int(input("Enter x: "))\\
y = int(input("Enter y: "))\\
z = int(input("Enter z: "))\\
print("Sum =", add(x, y, z))
}
{
\textbf{Function Contract:}

Input → তিনটি Number

Processing → Addition

Output → Sum
}

%================================================
% 58
%================================================

\exerciseTable
{৫৮. একই Function একাধিকবার Call}
{
১. Start।\\
২. Function Define করো।\\
৩. Function Call করো।\\
৪. আবার Function Call করো।\\
৫. End।
}
{
\linearflow
{Function Define}
{Function দুইবার Call}
{একই কাজ দুইবার Execute}
}
{
def hello():\\
\hspace*{4mm}print("Welcome")\\
\hspace*{4mm}print("Python")\\
hello()\\
hello()
}
{
একই Code আবার না লিখেও Function পুনরায় ব্যবহার করা যায়।
}

%================================================
% PART F
%================================================

\partheading{অংশ–F : গুরুত্বপূর্ণ Mathematical Program}

%================================================
% 59
%================================================

\exerciseTable
{৫৯. GCD / গ.সা.গু নির্ণয়}
{
১. Start।\\
২. \(a,b\) Input নাও।\\
৩. \(b\neq0\) পর্যন্ত remainder বের করো।\\
৪. \(a=b,\;b=remainder\)।\\
৫. শেষে \(a\) হলো GCD।\\
৬. End।
}
{
\linearflow
{a,b Input}
{Euclidean Remainder Process}
{GCD প্রদর্শন}
}
{
a = int(input("Enter a: "))\\
b = int(input("Enter b: "))\\
while b != 0:\\
\hspace*{4mm}r = a \% b\\
\hspace*{4mm}a = b\\
\hspace*{4mm}b = r\\
print("GCD =", a)
}
{
Euclidean Algorithm ব্যবহার করা হয়েছে।
}

%================================================
% 60
%================================================

\exerciseTable
{৬০. LCM / ল.সা.গু নির্ণয়}
{
১. Start।\\
২. \(a,b\) Input নাও।\\
৩. বড় Value থেকে Search শুরু করো।\\
৪. উভয় দ্বারা বিভাজ্য না হওয়া পর্যন্ত বাড়াও।\\
৫. LCM Print করো।\\
৬. End।
}
{
\linearflow
{a,b Input}
{Common Multiple Search}
{LCM প্রদর্শন}
}
{
a = int(input("Enter a: "))\\
b = int(input("Enter b: "))\\
lcm = max(a, b)\\
while True:\\
\hspace*{4mm}if lcm \% a == 0 and lcm \% b == 0:\\
\hspace*{8mm}break\\
\hspace*{4mm}lcm += 1\\
print("LCM =", lcm)
}
{
প্রথম Common Multiple পাওয়া গেলে Loop বন্ধ হয়।
}

%================================================
% 61
%================================================

\exerciseTable
{৬১. Leap Year নির্ণয়}
{
১. Start।\\
২. Year Input নাও।\\
৩. 400 দ্বারা বিভাজ্য হলে Leap Year।\\
৪. অথবা 4 দ্বারা বিভাজ্য কিন্তু 100 দ্বারা নয় হলে Leap Year।\\
৫. অন্যথায় Not Leap Year।\\
৬. End।
}
{
\linearflow
{Year Input}
{Leap Year Condition পরীক্ষা}
{Result প্রদর্শন}
}
{
year = int(input("Enter year: "))\\
if ((year \% 4 == 0 and year \% 100 != 0)\\
\hspace*{8mm}or year \% 400 == 0):\\
\hspace*{4mm}print("Leap Year")\\
else:\\
\hspace*{4mm}print("Not Leap Year")
}
{
Century Year-এর ক্ষেত্রে 400 দ্বারা বিভাজ্যতার Condition গুরুত্বপূর্ণ।
}

%================================================
% PART G
%================================================

\partheading{অংশ–G : Formula ভিত্তিক Program Bank}

\begin{rememberbox}

Formula-based Program-এর সাধারণ ধাপ—

\begin{center}

\textbf{Start}
\(\rightarrow\)
\textbf{Input}
\(\rightarrow\)
\textbf{Formula প্রয়োগ}
\(\rightarrow\)
\textbf{Result}
\(\rightarrow\)
\textbf{End}

\end{center}

প্রতিটি Problem-এর নিজস্ব Algorithm, Flowchart এবং Python Code
একই Table-এর মধ্যে দেওয়া হয়েছে।

\end{rememberbox}

%================================================
% 62
%================================================

\formulaexercise
{৬২}
{Celsius থেকে Fahrenheit}
{\(C\)}
{c = float(input("Celsius: "))}
{f = 9 * c / 5 + 32}
{f}
{\(F=\frac{9C}{5}+32\)}

%================================================
% 63
%================================================

\formulaexercise
{৬৩}
{Celsius থেকে Kelvin}
{\(C\)}
{c = float(input("Celsius: "))}
{k = c + 273.15}
{k}
{\(K=C+273.15\)}

%================================================
% 64
%================================================

\formulaexercise
{৬৪}
{Fahrenheit থেকে Celsius}
{\(F\)}
{f = float(input("Fahrenheit: "))}
{c = 5 * (f - 32) / 9}
{c}
{\(C=\frac{5(F-32)}{9}\)}

%================================================
% 65
%================================================

\formulaexercise
{৬৫}
{Fahrenheit থেকে Kelvin}
{\(F\)}
{f = float(input("Fahrenheit: "))}
{k = 5 * (f - 32) / 9 + 273.15}
{k}
{\(K=\frac{5(F-32)}{9}+273.15\)}

%================================================
% 66
%================================================

\formulaexercise
{৬৬}
{Kelvin থেকে Fahrenheit}
{\(K\)}
{k = float(input("Kelvin: "))}
{f = 9 * (k - 273.15) / 5 + 32}
{f}
{\(F=\frac{9(K-273.15)}{5}+32\)}

%================================================
% 67
%================================================

\formulaexercise
{৬৭}
{Kelvin থেকে Celsius}
{\(K\)}
{k = float(input("Kelvin: "))}
{c = k - 273.15}
{c}
{\(C=K-273.15\)}

%================================================
% 68
%================================================

\formulaexercise
{৬৮}
{Centimeter থেকে Inch}
{\(cm\)}
{cm = float(input("Centimeter: "))}
{inch = cm / 2.54}
{inch}
{\(inch=\frac{cm}{2.54}\)}

%================================================
% 69
%================================================

\formulaexercise
{৬৯}
{Inch থেকে Centimeter}
{\(inch\)}
{inch = float(input("Inch: "))}
{cm = inch * 2.54}
{cm}
{\(cm=inch\times2.54\)}

%================================================
% 70
%================================================

\formulaexercise
{৭০}
{Foot থেকে Meter}
{\(ft\)}
{ft = float(input("Foot: "))}
{m = ft / 3.28}
{m}
{\(m=\frac{ft}{3.28}\)}

%================================================
% 71
%================================================

\formulaexercise
{৭১}
{Meter থেকে Foot}
{\(m\)}
{m = float(input("Meter: "))}
{ft = m * 3.28}
{ft}
{\(ft=m\times3.28\)}

%================================================
% 72
%================================================

\formulaexercise
{৭২}
{বৃত্তের ক্ষেত্রফল}
{\(r\)}
{r = float(input("Radius: "))}
{area = 3.1416 * r * r}
{area}
{\(A=\pi r^2\)}

%================================================
% 73
%================================================

\formulaexercise
{৭৩}
{বৃত্তের পরিধি}
{\(r\)}
{r = float(input("Radius: "))}
{circumference = 2 * 3.1416 * r}
{circumference}
{\(C=2\pi r\)}

%================================================
% 74
%================================================

\formulaexercise
{৭৪}
{আয়তক্ষেত্রের ক্ষেত্রফল}
{\(l,w\)}
{l = float(input("Length: "))\\
w = float(input("Width: "))}
{area = l * w}
{area}
{\(A=lw\)}

%================================================
% 75
%================================================

\formulaexercise
{৭৫}
{আয়তক্ষেত্রের পরিসীমা}
{\(l,w\)}
{l = float(input("Length: "))\\
w = float(input("Width: "))}
{p = 2 * (l + w)}
{p}
{\(P=2(l+w)\)}

%================================================
% 76
%================================================

\formulaexercise
{৭৬}
{বর্গক্ষেত্রের ক্ষেত্রফল}
{\(a\)}
{a = float(input("Side: "))}
{area = a * a}
{area}
{\(A=a^2\)}

%================================================
% 77
%================================================

\formulaexercise
{৭৭}
{বর্গক্ষেত্রের পরিসীমা}
{\(a\)}
{a = float(input("Side: "))}
{p = 4 * a}
{p}
{\(P=4a\)}

%================================================
% 78
%================================================

\formulaexercise
{৭৮}
{সামান্তরিকের ক্ষেত্রফল}
{\(b,h\)}
{b = float(input("Base: "))\\
h = float(input("Height: "))}
{area = b * h}
{area}
{\(A=bh\)}

%================================================
% 79
%================================================

\formulaexercise
{৭৯}
{সামান্তরিকের পরিসীমা}
{\(a,b\)}
{a = float(input("Side a: "))\\
b = float(input("Side b: "))}
{p = 2 * (a + b)}
{p}
{\(P=2(a+b)\)}

%================================================
% 80
%================================================

\formulaexercise
{৮০}
{রম্বসের ক্ষেত্রফল}
{\(d_1,d_2\)}
{d1 = float(input("Diagonal 1: "))\\
d2 = float(input("Diagonal 2: "))}
{area = d1 * d2 / 2}
{area}
{\(A=\frac{d_1d_2}{2}\)}

%================================================
% 81
%================================================

\formulaexercise
{৮১}
{রম্বসের পরিসীমা}
{\(a\)}
{a = float(input("Side: "))}
{p = 4 * a}
{p}
{\(P=4a\)}

%================================================
% 82
%================================================

\formulaexercise
{৮২}
{ত্রিভুজের ক্ষেত্রফল}
{\(b,h\)}
{b = float(input("Base: "))\\
h = float(input("Height: "))}
{area = b * h / 2}
{area}
{\(A=\frac{bh}{2}\)}

%================================================
% 83
%================================================

\formulaexercise
{৮৩}
{ত্রিভুজের পরিসীমা}
{\(a,b,c\)}
{a = float(input("a: "))\\
b = float(input("b: "))\\
c = float(input("c: "))}
{p = a + b + c}
{p}
{\(P=a+b+c\)}

%================================================
% 84
%================================================

\formulaexercise
{৮৪}
{ত্রিভুজের অর্ধপরিসীমা}
{\(a,b,c\)}
{a = float(input("a: "))\\
b = float(input("b: "))\\
c = float(input("c: "))}
{s = (a + b + c) / 2}
{s}
{\(s=\frac{a+b+c}{2}\)}

%================================================
% 85
%================================================

\exerciseTable
{৮৫. Heron's Formula দিয়ে ত্রিভুজের ক্ষেত্রফল}
{
১. Start।\\
২. \(a,b,c\) Input নাও।\\
৩. \(s=(a+b+c)/2\)।\\
৪. \(A=\sqrt{s(s-a)(s-b)(s-c)}\)।\\
৫. Area প্রদর্শন করো।\\
৬. End।
}
{
\linearflow
{a,b,c Input}
{s ও Heron's Formula প্রয়োগ}
{area প্রদর্শন}
}
{
a = float(input("a: "))\\
b = float(input("b: "))\\
c = float(input("c: "))\\
s = (a + b + c) / 2\\
area = (s * (s-a) * (s-b) * (s-c)) ** 0.5\\
print("Area =", area)
}
{
প্রথমে Semi-perimeter নির্ণয় করা হয়।

\medskip

তারপর Heron's Formula প্রয়োগ করা হয়।
}

%================================================
% 86
%================================================

\formulaexercise
{৮৬}
{সমবাহু ত্রিভুজের ক্ষেত্রফল}
{\(a\)}
{a = float(input("Side: "))}
{area = (3 ** 0.5) * a * a / 4}
{area}
{\(A=\frac{\sqrt3\,a^2}{4}\)}

%================================================
% 87
%================================================

\formulaexercise
{৮৭}
{সমবাহু ত্রিভুজের পরিসীমা}
{\(a\)}
{a = float(input("Side: "))}
{p = 3 * a}
{p}
{\(P=3a\)}

%================================================
% 88
%================================================

\formulaexercise
{৮৮}
{সমদ্বিবাহু ত্রিভুজের ক্ষেত্রফল}
{\(a,b\)}
{a = float(input("Equal side: "))\\
b = float(input("Base: "))}
{area = b * ((4*a*a - b*b) ** 0.5) / 4}
{area}
{\(A=\frac{b\sqrt{4a^2-b^2}}{4}\)}

%================================================
% 89
%================================================

\formulaexercise
{৮৯}
{Trapezoid-এর ক্ষেত্রফল}
{\(a,b,h\)}
{a = float(input("Side a: "))\\
b = float(input("Side b: "))\\
h = float(input("Height: "))}
{area = (a + b) * h / 2}
{area}
{\(A=\frac{(a+b)h}{2}\)}

%================================================
% 90
%================================================

\formulaexercise
{৯০}
{Trapezoid-এর পরিসীমা}
{\(a,b,c,d\)}
{a = float(input("a: "))\\
b = float(input("b: "))\\
c = float(input("c: "))\\
d = float(input("d: "))}
{p = a + b + c + d}
{p}
{\(P=a+b+c+d\)}

%================================================
% 91
%================================================

\formulaexercise
{৯১}
{Cuboid-এর Surface Area}
{\(a,b,c\)}
{a = float(input("a: "))\\
b = float(input("b: "))\\
c = float(input("c: "))}
{area = 2 * (a*b + b*c + c*a)}
{area}
{\(A=2(ab+bc+ca)\)}

%================================================
% 92
%================================================

\formulaexercise
{৯২}
{Cuboid-এর Diagonal}
{\(a,b,c\)}
{a = float(input("a: "))\\
b = float(input("b: "))\\
c = float(input("c: "))}
{d = (a*a + b*b + c*c) ** 0.5}
{d}
{\(d=\sqrt{a^2+b^2+c^2}\)}

%================================================
% 93
%================================================

\formulaexercise
{৯৩}
{Cuboid-এর Volume}
{\(a,b,c\)}
{a = float(input("a: "))\\
b = float(input("b: "))\\
c = float(input("c: "))}
{v = a * b * c}
{v}
{\(V=abc\)}

%================================================
% 94
%================================================

\formulaexercise
{৯৪}
{Cube-এর Diagonal}
{\(a\)}
{a = float(input("Side: "))}
{d = (3 ** 0.5) * a}
{d}
{\(d=\sqrt3\,a\)}

%================================================
% 95
%================================================

\formulaexercise
{৯৫}
{Cube-এর Volume}
{\(a\)}
{a = float(input("Side: "))}
{v = a ** 3}
{v}
{\(V=a^3\)}

%================================================
% 96
%================================================

\formulaexercise
{৯৬}
{Cube-এর Surface Area}
{\(a\)}
{a = float(input("Side: "))}
{area = 6 * a * a}
{area}
{\(A=6a^2\)}

%================================================
% 97
%================================================

\formulaexercise
{৯৭}
{Cone-এর Slant Height}
{\(r,h\)}
{r = float(input("Radius: "))\\
h = float(input("Height: "))}
{l = (r*r + h*h) ** 0.5}
{l}
{\(l=\sqrt{r^2+h^2}\)}

%================================================
% 98
%================================================

\formulaexercise
{৯৮}
{Cone-এর Curved Surface Area}
{\(r,l\)}
{r = float(input("Radius: "))\\
l = float(input("Slant height: "))}
{area = 3.1416 * r * l}
{area}
{\(A=\pi rl\)}

%================================================
% 99
%================================================

\formulaexercise
{৯৯}
{Cone-এর Base Area}
{\(r\)}
{r = float(input("Radius: "))}
{area = 3.1416 * r * r}
{area}
{\(A=\pi r^2\)}

%================================================
% 100
%================================================

\formulaexercise
{১০০}
{Cone-এর Volume}
{\(r,h\)}
{r = float(input("Radius: "))\\
h = float(input("Height: "))}
{v = 3.1416 * r * r * h / 3}
{v}
{\(V=\frac{\pi r^2h}{3}\)}

%================================================
% 101
%================================================

\formulaexercise
{১০১}
{Square Root নির্ণয়}
{\(n\)}
{n = float(input("Number: "))}
{root = n ** 0.5}
{root}
{\(root=\sqrt n\)}

%================================================
% 102
%================================================

\exerciseTable
{১০২. Absolute Value নির্ণয়}
{
১. Start।\\
২. Number Input নাও।\\
৩. Absolute Value নির্ণয় করো।\\
৪. Result প্রদর্শন করো।\\
৫. End।
}
{
\linearflow
{n Input}
{value=abs(n)}
{value প্রদর্শন}
}
{
n = float(input("Number: "))\\
value = abs(n)\\
print("Absolute =", value)
}
{
Negative Input হলেও Result non-negative হবে।
}

%================================================
% 103
%================================================

\formulaexercise
{১০৩}
{Power নির্ণয়}
{\(a,b\)}
{a = float(input("Base: "))\\
b = float(input("Power: "))}
{p = a ** b}
{p}
{\(p=a^b\)}

%================================================
% Revision Map
%================================================

\chFiveInput{topics/47_algorithm_flowchart_exercises/python_extensions.tex}

\partheading{Exercise Revision Map}

\begin{tabularx}{\linewidth}{
|>{\bfseries\RaggedRight\arraybackslash}p{0.29\linewidth}
|Y
|Y|
}

\hline
\rowcolor{softBlue}

\textbf{Problem Type}
&
\textbf{মূল Logic}
&
\textbf{Python Concept}
\\
\hline

Sum / Average
&
Accumulator
&
Operator + Loop
\\
\hline

Largest / Smallest
&
Comparison
&
if-elif-else
\\
\hline

Even / Odd
&
Remainder
&
\textenglish{\ttfamily \%}
\\
\hline

Grade / Range
&
Multiple Condition
&
if-elif-else
\\
\hline

Series
&
Repeated Term Generation
&
for / while
\\
\hline

Factorial
&
Repeated Multiplication
&
Loop
\\
\hline

Prime
&
Divisor Test
&
Nested Loop + break
\\
\hline

List
&
Traversal
&
for
\\
\hline

Search
&
Sequential Comparison
&
Loop + break
\\
\hline

Function
&
Reusable Logic
&
def + return
\\
\hline

GCD
&
Repeated Remainder
&
while
\\
\hline

LCM
&
Common Multiple Search
&
while
\\
\hline

Formula Program
&
Input → Formula → Output
&
Arithmetic Operator
\\
\hline

\end{tabularx}

%================================================
% Common Mistakes
%================================================

\partheading{Exercise-এ সাধারণ ভুল}

\begin{rememberbox}

\begin{itemize}

\item Algorithm-এ Start ও End বাদ দেওয়া।

\item Flowchart-এ Input/Output এবং Processing-এর পার্থক্য না রাখা।

\item Decision Flowchart-এ Yes/No Branch না দেখানো।

\item Python Code-এ Indentation ভুল করা।

\item
\textenglish{\ttfamily input()}
থেকে পাওয়া Data প্রয়োজনমতো
\textenglish{\ttfamily int()}
বা
\textenglish{\ttfamily float()}
-এ Convert না করা।

\item
\textenglish{\ttfamily =}
ও
\textenglish{\ttfamily ==}
গুলিয়ে ফেলা।

\item
\textenglish{\ttfamily range()}
-এর Stop Value Excluded ভুলে যাওয়া।

\item Sum-এর Initial Value 0 এবং Factorial-এর Initial Value 1—এটি ভুল করা।

\item
\textenglish{\ttfamily while}
Loop-এ Update না দেওয়া।

\item Prime Number পরীক্ষায় 0 ও 1-কে Prime মনে করা।

\item Formula-তে Bracket ভুল দেওয়া।

\item Python-এ Power-এর জন্য
\textenglish{\ttfamily **}
ব্যবহার করতে হয়।

\end{itemize}

\end{rememberbox}

%================================================
% Final Reminder
%================================================

\begin{rememberbox}

\textbf{\color{deepOrange}মনে রাখবে}

\begin{center}

\textbf{Algorithm}
\(\longrightarrow\)
Problem সমাধানের ধাপ

\medskip

\textbf{Flowchart}
\(\longrightarrow\)
Algorithm-এর চিত্ররূপ

\medskip

\textbf{Python Program}
\(\longrightarrow\)
Algorithm-এর বাস্তব Code

\medskip

\textbf{Output}
\(\longrightarrow\)
Program চালানোর ফলাফল

\end{center}

\end{rememberbox}

\small

%================================================
% TITLE
%================================================
